\documentclass[11pt,a4paper]{article}
\usepackage[T1]{fontenc}
\usepackage[utf8]{inputenc}
\usepackage[spanish,es-noshorthands]{babel}
\usepackage{amsmath,amssymb}
\usepackage{graphicx,booktabs,hyperref,listings,geometry}
\geometry{margin=2.4cm}
\lstset{basicstyle=\ttfamily\small,breaklines=true,columns=fullflexible}
\hypersetup{colorlinks=true,linkcolor=black,urlcolor=blue}

\title{Documentación del software \texttt{fibonacci\_tmm}}
\author{Emanuel Orozco Gallego}
\date{31 de agosto de 2026}

\setlength{\emergencystretch}{3em}
\begin{document}
\maketitle

\begin{abstract}
Manual independiente de la biblioteca Python que reproduce
Reyes-Gómez \textit{et al.}, Phys.\ Rev.\ B \textbf{81}, 153101 (2010),
con el método de matriz de transferencia.
\end{abstract}

\tableofcontents
\newpage

\section{Arquitectura}

El paquete vive en \texttt{src/fibonacci\_tmm/}.

\begin{center}
\begin{tabular}{ll}
\toprule
Módulo & Responsabilidad \\
\midrule
\texttt{fibonacci.py} & $S_m$, $F_m$, $L_m$, inflación \\
\texttt{materials.py} & Aire y Drude \\
\texttt{electromagnetics.py} & $n$, $q$, $Q$, TE/TM \\
\texttt{transfer\_matrix.py} & $M_j$, $T_m$, $R_m$ \\
\texttt{dispersion.py} & $\cos(kL)=R$, bandas \\
\texttt{plasmon\_modes.py} & recuento de subbandas \\
\texttt{bandwidth.py} & intervalos $\nu_{\min}(\theta)$--$\nu_{\max}(\theta)$ (bordes refinados) \\
\texttt{params.py} & YAML $\to$ SI \\
\texttt{plotting.py} & estilo de figuras \\
\texttt{provenance.py} & metadatos de corrida \\
\bottomrule
\end{tabular}
\end{center}

Los parámetros científicos no están en los scripts: están en \texttt{configs/figure\_0N.yaml}.

\section{Instalación}

\begin{lstlisting}
python3 -m venv .venv
source .venv/bin/activate
pip install -r requirements.txt
\end{lstlisting}

Python $\ge$ 3.10. Dependencias: NumPy, SciPy, Matplotlib, PyYAML, pytest.

\section{Uso}

\begin{lstlisting}
.venv/bin/pytest
.venv/bin/python scripts/tmm/figure_01.py      # una figura, TMM
.venv/bin/python scripts/tmm/run_all.py        # Figs. 1-6, TMM
.venv/bin/python scripts/pwe/run_all.py        # Figs. 1-6, ondas planas
.venv/bin/python scripts/comparison/run_all.py # TMM vs PWE
make paper pwe-docs
\end{lstlisting}

Para un orden nuevo de Fibonacci, cambie \texttt{fibonacci\_orders} en el YAML.
Para TM, ponga \texttt{polarization: TM}.
Para otro Drude, edite \texttt{omega\_e\_over\_2pi\_ghz} y \texttt{omega\_m\_over\_2pi\_ghz}.

\section{API principal}

\texttt{sequence(m)}, \texttt{paper\_fibonacci(m)}, \texttt{build\_cell(m,a,b)}.

{\raggedright
\texttt{scan\_dispersion(spec, m, theta, nu\_ghz, polarization)} (TMM; alias de
\texttt{scan\_dispersion\_tmm}).

\texttt{scan\_dispersion\_pwe(spec, m, theta, nu\_ghz, polarization,
harmonics\_per\_layer=16)}: la misma semitraza con ondas planas; devuelve el
mismo \texttt{DispersionScan}.\par}

\texttt{compare\_scans(ref, cand)}, \texttt{locate\_bands},
\texttt{validate\_bands}: comparación y localización de subbandas.

El paquete es \texttt{fibonacci\_photonics} (subpaquetes \texttt{core},
\texttt{tmm}, \texttt{pwe}, \texttt{analysis}); \texttt{fibonacci\_tmm}
sigue disponible como alias.

\texttt{semitrace\_by\_recurrence} y \texttt{semitrace\_by\_product} deben coincidir.

\texttt{detect\_plasmon\_modes(scan, nu\_min\_ghz=..., nu\_max\_ghz=...)}.

\texttt{load\_spec(path)} lee un YAML.

\section{Modelo físico}

Aire $\varepsilon=\mu=1$. Drude $\varepsilon=1-\omega_e^2/\omega^2$,
$\mu=1-\omega_m^2/\omega^2$. Frecuencia graficada: $\nu=\omega/2\pi$ en GHz.

\section{TMM}

Matriz de capa unimodular con $\chi=\mu$ (TE) o $\varepsilon$ (TM).
Producto $T=M_N\cdots M_1$. Recurrencia $R_m=2R_{m-1}R_{m-2}-R_{m-3}$.
$q=(\omega/c)n_A\sin\theta$ (el paper omite $\omega/c$).

\section{Fibonacci}

$S_0=B$, $S_1=A$, $S_m=S_{m-1}S_{m-2}$, $F_0=F_1=1$,
$N_B=F_{m-2}$, $L_m=F_{m-1}a+F_{m-2}b$.

\section{Dispersión}

Banda si $\lvert R\rvert\le 1+10^{-12}$. $kL_m/\pi=\arccos(R)/\pi$.
Los puntos con $\lvert R\rvert>1$ se dejan como NaN, no se recortan en silencio.

\section{Tests}

\begin{itemize}
\item \texttt{core/test\_fibonacci.py}: secuencias, inflación, $N_B$, $L_m$.
\item \texttt{core/test\_materials.py}: ceros de Drude, $\nu_e,\nu_m$.
\item \texttt{core/test\_params.py}: YAML de las figuras.
\item \texttt{tmm/test\_transfer\_matrix.py}: $\det T=1$, (7)--(9), recurrencia vs producto, TE$=$TM a $\theta=0$.
\item \texttt{tmm/test\_dispersion.py}: máscara de banda, gap $\langle n\rangle=0$.
\item \texttt{pwe/test\_fourier.py}, \texttt{pwe/test\_solver.py}: Fourier de la celda, benchmark del libro, $R_{\rm PWE}\approx R_{\rm TMM}$.
\item \texttt{analysis/test\_validation.py}: $N_{\mathrm{modos}}=F_{m-2}$.
\item \texttt{analysis/test\_comparison.py}: métricas y localización de bordes.
\item \texttt{compat/test\_legacy\_imports.py}: alias \texttt{fibonacci\_tmm}.
\end{itemize}

\section{Reproducción de figuras}

Cada \texttt{results/tmm/figure\_0N/} contiene \texttt{data/}, \texttt{output/} y \texttt{README.md}
(\texttt{figures/} es un enlace a \texttt{results/tmm/}). Las versiones PWE están en
\texttt{results/pwe/figure\_0N/} y la comparación en \texttt{results/comparison/}.
Las imágenes finales son PDF, PNG (300 dpi) y SVG.

\section{Solución de problemas}

\begin{itemize}
\item \texttt{ZeroDivisionError} en $\omega=0$: el Drude diverge; no incluya el origen en el barrido.
\item $\mu_B=0$ o $\varepsilon_B=0$: polo; el escáner enmascara una vecindad relativa $10^{-6}$.
\item $\lvert R\rvert$ siempre $\le 1$ a $\theta=0$ si $\omega_e=\omega_m$: no es un bug, es adaptación de impedancia.
\item Fig.~6: el original grafica $\nu(\theta)$ relleno, no $\Delta\nu$;
el caption del PRB está incompleto (parámetros inferidos de las Figs.~3--5).
\item ImportError: añada \texttt{src} al \texttt{PYTHONPATH} o use los scripts de \texttt{scripts/}.
\end{itemize}

\section{Extensiones}

Nuevo orden: YAML. Nueva polarización: \texttt{Polarization.TM}. Nuevo material: implemente \texttt{epsilon\_mu(omega)} con la misma interfaz que \texttt{DrudeMetamaterial}.

\end{document}
